

Il est
évidemment nécessaire pour cela que @@M0@@ soit caténaire, et suffisant que @@M1@@ soit biéquidimensionnel (EGA @@M2@@ 14.3.3). Cependant il existe des schémas noethériens de dimension finie, caténaires et équidimensionnels, n'ayant pas de telles fonctions @@M3@@. Pour le voir on peut modifier l'exemple (EGA IV 5.6.11) de la façon suivante : prenons deux exemplaires @@M4@@ et @@M5@@ du spectre de l'anneau @@M6@@ de loc. cit. et désignons par @@M7@@ et @@M8@@ (@@M9@@) les points fermés de @@M10@@ (@@M11@@) correspondant aux idéaux maximaux @@M12@@ resp. @@M13@@ de @@M14@@ de hauteur @@M15@@ resp. @@M16@@. Rappelons que les corps résiduels de @@M17@@ et @@M18@@ sont isomorphes. On peut donc construire par recollement de @@M19@@ et @@M20@@ un schéma @@M21@@ en identifiant @@M22@@ avec @@M23@@ et @@M24@@ avec @@M25@@. On peut alors identifier @@M26@@ et @@M27@@ avec les deux composantes irréductibles de @@M28@@, et alors dans @@M29@@, @@M30@@ est de dimension zéro. On voit alors que @@M31@@ est un schéma noethérien, caténaire et équidimensionnel (de dimension @@M32@@). Si (ii bis) était vérifié sur @@M33@@, on aurait
@@M34@@
et d'autre part
@@M35@@
donc on trouve une contradiction.

